
\section*{1. Énoncé du théorème principal et applications}

\textbf{Théorème 1.1.} \uline{Soient} @@M0@@ \uline{un schéma noethérien,} @@M1@@ \uline{et} @@M2@@ \uline{deux} @@M3@@\uline{-schémas propres,} @@M4@@ \uline{un morphisme. On suppose} @@M5@@ intègre \uline{et} @@M6@@ surjectif. \uline{Alors il existe un ouvert non-vide} @@M7@@ \uline{de} @@M8@@ \uline{tel que le morphisme}
@@M9@@
\uline{entre les foncteurs de Picard (localisés pour la topologie} f.p.p.f) \uline{soit représentable par un morphisme quasi-affine de présentation finie.}

\textbf{Corollaire 1.2.} \uline{Soit} @@M10@@ \uline{un schéma propre sur un schéma noethérien} @@M11@@\uline{. On suppose} @@M12@@ intègre. \uline{Alors il existe un ouvert non vide} @@M13@@ \uline{de} @@M14@@ \uline{tel que le foncteur de Picard} @@M15@@ \uline{soit représentable par un} @@M16@@\uline{-schéma en groupes somme d'ouverts quasi-projectifs de type fini, donc a fortiori séparé sur} @@M17@@.

En effet rappelons le fait suivant (FGA N\textsuperscript{o}232)~: soient @@M18@@ un schéma, @@M19@@ un @@M20@@-schéma projectif, de présentation finie et intègre sur @@M21@@ (c'est-à-dire plat à fibres géométriquement intègres). Alors le foncteur de Picard est représenté par un @@M22@@-schéma en groupes @@M23@@, réunion croissante d'ouverts quasi-projectifs de présentation finie.

La considération des «~polynômes de Hilbert~» permet ici de décomposer @@M24@@ en somme d'ouverts @@M25@@, @@M26@@. Dans l'exposé suivant on démontre sous l'hypothèse supplémentaire @@M27@@ noethérien que les @@M28@@ sont de type fini sur @@M29@@. On sait alors que 1.2 est vrai dans le cas où @@M30@@ est en plus projectif et intègre sur @@M31@@; on va en déduire le cas général en utilisant le

\textbf{Lemme 1.3.} \uline{Soit} @@M32@@ \uline{de type fini sur} @@M33@@ \uline{noethérien intègre. Alors il existe un ouvert non vide} @@M34@@ \uline{de} @@M35@@\uline{, un schéma intègre} @@M36@@ \uline{et un morphisme fini et}
